\documentclass[11pt,a4paper]{article}
\usepackage{operationsform}
\begin{document}
\OPStitle{企画書}{TPL-DOC-030}{0.2.0}
\OPSsection{1}{基本情報}
\OPSblanklong{案件・文書の識別子}
\OPSblankshort{作成者と連絡先}
\OPSblankshort{作成日と対象版}
\OPSsection{2}{目的と対象読者}
\OPSblanklong{企画の背景・公共的意義}
\OPSblanklong{編集目的・主要な問い}
\OPSblanklong{想定読者と前提知識}
\OPSblanklong{記事の中心的な範囲}
\OPSsection{3}{取材と制作計画}
\OPSblanklong{根拠資料と未確認事項}
\OPSblanklong{取材先と確認事項}
\OPSblanklong{予定構成・媒体}
\par\Needspace{4\baselineskip}\noindent\textbf{工程・役割・確認経路}\par\vspace{1mm}
{\fontsize{8.5}{13}\selectfont
\begin{tabularx}{\linewidth}{@{}XXXX@{}}
\textbf{対象・項目} & \textbf{内容・参照} & \textbf{担当・状態} & \textbf{備考}\\[2mm]
\strut & \strut & \strut & \strut\\[5mm]
\strut & \strut & \strut & \strut\\[5mm]
\end{tabularx}
}\vspace{2mm}\par
\OPSsection{4}{確認と履歴}
\par\Needspace{4\baselineskip}\noindent\textbf{参照資料と根拠ID}\par\vspace{1mm}
{\fontsize{8.5}{13}\selectfont
\begin{tabularx}{\linewidth}{@{}XXXX@{}}
\textbf{資料ID} & \textbf{出典・該当箇所} & \textbf{版・日付} & \textbf{確認状態}\\[2mm]
\strut & \strut & \strut & \strut\\[5mm]
\strut & \strut & \strut & \strut\\[5mm]
\end{tabularx}
}\vspace{2mm}\par
\OPSblanklong{未確認事項・例外}
\par\Needspace{4\baselineskip}\noindent\textbf{点検・改訂の記録}\par\vspace{1mm}
{\fontsize{8.5}{13}\selectfont
\begin{tabularx}{\linewidth}{@{}XXXX@{}}
\textbf{日時・版} & \textbf{確認者} & \textbf{判断・理由} & \textbf{変更点}\\[2mm]
\strut & \strut & \strut & \strut\\[5mm]
\strut & \strut & \strut & \strut\\[5mm]
\end{tabularx}
}\vspace{2mm}\par
\end{document}
